% !TEX root = 00_instructivo.tex
\chapter{Transductores acústicos: Monitoreo Submarino}

\section*{Objetivos}
\begin{itemize}
    \item Construir un hidrófono básico usando un sensor piezoeléctrico impermeable.
    \item Amplificar la señal del piezoeléctrico con un módulo MAX4466.
    \item Clasificar niveles de presión sonora subacuática (decibeles relativos) mediante un algoritmo de pico a pico.
    \item Activar indicadores visuales (LEDs) según la intensidad del sonido detectado.
\end{itemize}

\section*{Preguntas previas}
\begin{enumerate}
    \item ¿Qué es el efecto piezoeléctrico directo y cómo se utiliza para detectar sonido?
    \item ¿Por qué un piezoeléctrico sumergido en agua actúa como un hidrófono?
    \item ¿Qué función cumple el amplificador MAX4466? ¿Por qué es necesario para un Arduino de 5V?
    \item Investigue: ¿Qué rango de frecuencias audibles puede detectar un cerámico PZT típico?
\end{enumerate}

\section*{Materiales y equipo}
\begin{itemize}
    \item Arduino Uno.
    \item Sensor piezoeléctrico con revestimiento impermeable (Keyestudio o similar).
    \item Módulo amplificador MAX4466 (ganancia ajustable).
    \item Recipiente transparente con agua (pecera o cubo).
    \item LEDs (azul y rojo) y resistencias de 220 $\Omega$.
    \item Mini protoboard.
    \item Cables Dupont.
    \item Fuente de sonido subacuática (puede ser un motor pequeño, un diapasón, o altavoz sumergible).
    \item Opcional: Motores de vibración (de celulares) para simular sonidos de animales acuáticos.
\end{itemize}

> Mía:
\section*{Procedimiento}
\begin{enumerate}
    \item \textbf{Preparación del sensor:}
    \begin{enumerate}
        \item Verifique que el sensor piezoeléctrico esté completamente impermeabilizado (si no, cúbralo con resina epóxica o silicona).
        \item Conecte el piezoeléctrico a la entrada del amplificador MAX4466 (IN y GND).
    \end{enumerate}
    \item \textbf{Conexión eléctrica:}
    \begin{enumerate}
        \item Conecte el MAX4466 (VCC a 3.3V, GND a GND, OUT a A0).
        \item Conecte los LEDs: ánodo largo a D6 (azul) y D7 (rojo) con resistencias de 220 $\Omega$ a GND.
        \item Cargue el código de la Sección 5.1 de la guía.
    \end{enumerate}
    \item \textbf{Calibración en aire y agua:}
    \begin{enumerate}
        \item En aire, ajuste la ganancia del MAX4466 (potenciómetro) hasta que el valor pico a pico en reposo sea alrededor de 200-300.
        \item Sumerja el sensor en el agua (solo la cápsula piezo, no el circuito).
        \item Observe el monitor serial. Ajuste \texttt{NIV\_NORM} y \texttt{NIV\_ALTO} según el ruido ambiental.
    \end{enumerate}
    \item \textbf{Pruebas de sonido:}
    \begin{enumerate}
        \item Genere sonidos suaves (golpear suavemente el borde del recipiente). Observe el LED azul.
        \item Genere sonidos fuertes (golpe fuerte, motor vibrando dentro del agua). Observe el LED rojo.
        \item Registre los valores pico a pico para cada tipo de sonido.
    \end{enumerate}
\end{enumerate}

\section*{Esquemático}
\begin{figure}[htbp]
\centering
\begin{circuitikz}[scale=0.8, transform shape]
\draw (0,0) node[above]{Arduino Uno};
\draw (1,0) node[ground]{GND};
\draw (1,0.5) node[anchor=south]{3.3V / 5V};
% MAX4466
\draw (0.5,1.5) node[anchor=east]{A0} -- (2.5,1.5);
\node at (3,1.5) {MAX4466 (OUT)};
% LEDs
\draw (0.5,2.5) node[anchor=east]{D6} -- (2.5,2.5);
\node at (3,2.5) {LED Azul};
\draw (0.5,3) node[anchor=east]{D7} -- (2.5,3);
\node at (3,3) {LED Rojo};
% Sensor piezo
\draw (2.5,1) node[anchor=west]{Piezo (IN)} -- (1,1);
\draw (1,1) to[short] (1,1.5);
\draw (2.5,0.5) node[anchor=west]{GND piezo} -- (1,0.5);
\end{circuitikz}
\caption{Conexión del hidrófono piezoeléctrico con amplificador MAX4466 y LEDs indicadores.}
\end{figure}

\section*{Preguntas de análisis}
\begin{enumerate}
    \item ¿Cómo se relaciona la amplitud pico a pico medida con la intensidad del sonido (decibeles)? ¿Es una relación lineal?
    \item ¿Qué diferencias observa en la sensibilidad del sensor dentro y fuera del agua? ¿Por qué?
    \item ¿Qué limitaciones tiene este hidrófono casero comparado con un hidrófono profesional?
    \item Explique por qué se usa una ventana de tiempo (\texttt{WINDOW}) para calcular el pico a pico en lugar de una muestra instantánea.
    \item Proponga una modificación al código para estimar la frecuencia del sonido (dominio del tiempo) y clasificar diferentes tipos de fuentes.
\end{enumerate}

\section*{Bibliografía}
\begin{itemize}
    \item Hoja de datos MAX4466. Maxim Integrated.
    \item “Piezoelectric sensors” – Measurement Specialties, Inc.
\end{itemize}
